\section{Kapitel 2}
\subsection{Bibeltext}

    \spa{1}\hh{1}Als solche, \\
    \spa{1}die also \verbN[b, p]{abgelegt haben} alle Schlechtigkeit \\
        \spa{2}und allen Trug \\
        \spa{2}und [alle] Heucheleien und Beneidungen \\
        \spa{2}und alle üblen Nachreden, \\
    \spa{1}\hh{2}\verbI[b, p]{seid} begierig wie neugeborene \person[b, p]{Kinder} \\
        nach der unverfälschten Milch des Wortes, \\
    \spa{1}damit ihr durch sie \verbI[b, p]{wachst}, \\
    \spa{1}\hh{3}wenn ihr wirklich \verbN[b, p]{schmecktet}, \\
        \spa{2}dass der \person[b, p]{Herr} freundlich ist, \\
        \spa{2}\hh{4}zu dem \verbN[b, p]{hinkommend}, \\
            \spa{3}einem lebenden Stein, \\
            \spa{3}von \person[b, p]{Menschen} \verbN[b, p]{abgelehnt} und \verbN[b, p]{verworfen}, \\
            \spa{3}ja, aber bei \person[b, p]{Gott} erwählt, \\
            \spa{3}kostbar, \\
        \spa{2}\hh{5}auch ihr selbst als lebende Steine \verbP[b, p]{gebaut werdet}: \\
            \spa{3}ein geistliches Haus, \\
            \spa{3}eine heilige Priesterschaft, \\
            \spa{3}um darzubringen geistliche Opfer, \\
            \spa{3}die \person[b, p]{Gott} angenehm sind durch \person[b, p]{Jesus Christus}, \\
        \spa{2}\hh{6}weswegen in der Schrift enthalten \verbN[b, p]{ist}: \\
            \spa{3}„Siehe, ich \verbN[b, p]{lege} in \ort[b, p]{Zion} einen Eckstein, \\
            \spa{3}einen \verbN[b, p]{erwählten}, kostbaren, \\
            \spa{3}und der, der an ihn \verbN[b, p]{glaubt}, \\
            \spa{3}\verbN[b, p]{wird} keineswegs \verbN[b, p]{zuschanden werden}.“ {Jes 28,16} \\

    \spa{1}\hh{7}Euch also, den \person[b, p]{Glaubenden}, \\
        \spa{2}[\verbN[b, p]{ist} er] die Kostbarkeit. \\
        \spa{2}Aber den im Unglauben Ungehorsamen \verbN[b, p]{wurde} dieser Stein, \\
        \spa{2}den die \person[b, p]{Bauenden} \verbN[b, p]{verwarfen}, zu einem Haupt[stein] der Ecke, \\
        \spa{2}\hh{8}und ein Stein des Stolperns [\verbN[b, p]{wurde} er] \\
        \spa{2}und ein Fels des Ärgernisses [denen], \\
            \spa{3}die als im Unglauben Ungehorsame sich am Wort \verbN[b, p]{stoßen}, \\
            \spa{3}\verbN[b, p]{stolpern} und \verbN[b, p]{fallen}, \\
            \spa{3}wozu sie auch \verbP[b, p]{gesetzt wurden}. \\

    \spa{1}\hh{9}Aber ihr \verbN[b, p]{seid} ein erwähltes \person[b, p]{Geschlecht}, \\
        \spa{2}eine königliche \person[b, p]{Priesterschaft}, \\
        \spa{2}ein heiliges \person[b, p]{Volk}, \\
        \spa{2}eine \person[b, p]{Volksschar} zu einem erworbenen Eigentum, \\
        \spa{2}auf dass ihr kund \verbN[b, p]{werden} lassen \verbN[b, p]{solltet} \\
        \spa{2}die lobenswerten Wesenszüge dessen, \\
            \spa{3}der euch aus der Finsternis in \verbN[b, p]{sein} wunderbares Licht \verbN[b, p]{rief}, \\
            \spa{3}\hh{10}die ihr einst „Nicht-Volk“ \verbN[b, p]{wart}, \\
            \spa{3}aber nun \person[b, p]{Gottes Volk} \verbN[b, p]{seid}, \\
                    \spa{4}die ihr „nicht Barmherzigkeit“ {Hos 1,8.9} \verbN[b, p]{empfangen hattet}, \\
                    \spa{4}nun aber Barmherzigkeit \verbN[b, p]{empfingt}. \\

    \spa{1}\hh{11}Geliebte, \\
        \spa{2}ich \verbN[b, p]{rufe} euch auf als \person[b, p]{‘Ausländer’} \\
        \spa{2}und sich vorübergehend aufhaltende \person[b, p]{Fremde}: \\
            \spa{3}\verbI[b, p]{Enthaltet} euch stets der fleischlichen Lüste – \\
            \spa{3}sie kämpfen gegen die Seele – \\
            \spa{3}\hh{12}und \verbI[b, p]{habt} eine edle Lebensführung \\
            \spa{3}unter denen von den \person[b, p]{Völkern}, \\
                \spa{4}damit sie da, \\
                \spa{4}wo sie gegen euch \verbN[b, p]{reden} wie gegen \person[b, p]{Übeltäter}, \\
                \spa{4}aufgrund der edlen Werke, \\
                \spa{4}die sie \verbN[b, p]{gesehen haben}, \\
                \spa{4}\person[b, p]{Gott} verherrlichen am Tag der Heimsuchung. \\
            \spa{3}\hh{13}\verbI[b, p]{Seid} also untergeordnet \\
                \spa{4}aller menschlichen Einrichtung des \person[b, p]{Herrn} wegen, \\
                \spa{4}\hh{14}\verbI[b, p]{sei} es dem \person[b, p]{König} als dem \person[b, p]{Übergeordneten} \\
                \spa{4}oder den \person[b, p]{Statthaltern} als denen, \\
                    \spa{5}die durch ihn \verbP[b, p]{geschickt wurden} zum Rechtsvollzug über \person[b, p]{Übeltäter}, \\
                \spa{4}aber zum Lob derer, \\
                    \spa{5}die Gutes \verbN[b, p]{tun}, \\
                \spa{4}\hh{15}weil es so \person[b, p]{Gottes} Wille \verbN[b, p]{ist}, \\
                    \spa{5}[durch] Gutestun die Unkenntnis \\
                    \spa{5}der törichten \person[b, p]{Menschen} zum Verstummen zu \verbN[b, p]{bringen} – \\
                \spa{4}\hh{16}als \person[b, p]{Freie} und nicht als solche, \\
                    die die Freiheit zum Deckmantel der Bosheit \verbN[b, p]{haben}, \\
                \spa{4}sondern als leibeigene \person[b, p]{Knechte Gottes}; \\
                    \spa{5}\hh{17}\verbI[b, p]{ehrt} alle; \\
                    \spa{5}\verbI[b, p]{liebt} die Bruderschaft; \\
                    \spa{5}\verbI[b, p]{fürchtet} \person[b, p]{Gott}; \\
                    \spa{5}\verbI[b, p]{ehrt} den \person[b, p]{König}; \\

        \spa{2}\hh{18}die \person[b, p]{Hausknechte} in aller Furcht sich unterordnend \\
            \spa{3}den [über sie] \verbN[b, p]{verfügenden} \person[b, p]{Herren}, \\
            \spa{3}nicht allein den guten und milden, \\
                \spa{4}sondern auch den krummen, \\
            \spa{3}\hh{19}denn das \verbN[b, p]{ist} Gnade, \\
            \spa{3}wenn jemand wegen eines an \person[b, p]{Gott} gebundenen Gewissens \\
                \spa{4}Betrübnisse und Verletzungen \verbN[b, p]{erträgt} \\
                \spa{4}und dabei ungerechterweise \verbN[b, p]{leidet}; \\
            \spa{3}\hh{20}denn was für ein Ruf \verbN[b, p]{ist} das, \\
                \spa{4}wenn ihr sündigt und [deswegen] mit Fäusten \verbP[b, p]{geschlagen werdet} \\
                \spa{4}und es mit Ausdauer erdulden \verbN[b, p]{werdet}? \\
            \spa{3}Wenn ihr jedoch Gutes \verbN[b, p]{tut} und \verbN[b, p]{leidet}, \\
                \spa{4}und es mit Ausdauer erdulden \verbN[b, p]{werdet}, \\
                \spa{4}das \verbN[b, p]{ist} Gnade bei \person[b, p]{Gott}, \\
            \spa{3}\hh{21}denn hierzu \verbP[b, p]{wurdet ihr gerufen}, \\
                \spa{4}weil auch \person[b, p]{Christus} für uns \verbN[b, p]{litt}, \\
                    \spa{5}wobei er euch ein vorgezeichnetes Muster \verbN[b, p]{hinterließ}, \\
                \spa{4}damit ihr auf seinen Fußspuren \verbN[b, p]{folgen möchtet}, \\
                    \spa{5}\hh{22}der keine Sünde \verbN[b, p]{tat}, \\
                    \spa{5}noch \verbP[b, p]{wurde} Trug in seinem Munde \verbP[b, p]{gefunden}, \\
                    \spa{5}\hh{23}der, als er geschmäht \verbP[b, p]{wurde}, \\
                    \spa{5}nicht \verbP[b, p]{wiederschmähte}, \\
                    \spa{5}als er \verbN[b, p]{litt}, \\
                    \spa{5}nicht \verbN[b, p]{drohte}, \\
                    \spa{5}es aber dem \verbN[b, p]{übergab}, \\
                    \spa{5}der in Gerechtigkeit \verbN[b, p]{richtet}, \\
                    \spa{5}\hh{24}der in seinem Leibe unsere Sünden selber auf das Holz \verbN[b, p]{hinauftrug}, \\
                \spa{4}damit wir, \\
                    \spa{5}den Sünden \verbN[b, p]{entledigt}, \\
                    \spa{5}der Gerechtigkeit leben \verbN[b, p]{möchten}, \\
                    \spa{5}durch dessen Wunde ihr \verbP[b, p]{geheilt wurdet}; \\
                \spa{4}\hh{25}denn ihr wart wie irrende Schafe; \\
                    \spa{5}ihr \verbN[b, p]{seid} jedoch nun \verbN[b, p]{zurückgekehrt} \\
                    \spa{5}zu dem \person[b, p]{Hirten} und \person[b, p]{Aufseher} eurer Seelen; \\
